\subsection{معرفی و انگیزه}
این ماژول در چارچوب پژوهش حاضر پیاده‌سازی و ارزیابی شده و فاقد مقاله مستقل جداگانه در ادبیات تحقیق است.

فشرده‌سازی مدل‌های یادگیری عمیق معمولاً در قالب دو فرآیند مجزا شامل هرس کردن وزن‌ها \lr{(Pruning)} و کوانتیزاسیون مقادیر \lr{(Quantization)} دنبال می‌شود. اجرای ترتیبی و گسسته این دو گام اغلب منجر به انباشت خطای تقریب و افت فاحش عملکرد مدل در نرخ‌های فشرده‌سازی بالا می‌گردد؛ چرا که هر گام بدون اطلاع از محدودیت‌های گام دیگر عمل می‌کند.

ماژول فشرده‌ساز \lr{ADMMCompressor} با به کارگیری روش جهت‌های متناوب ضرایب چندگانه لاگرانژ \lr{(Alternating Direction Method of Multipliers - ADMM)}، چارچوبی یکپارچه و بهینه‌سازی‌شده برای فشرده‌سازی توأم وزن‌ها ارائه می‌دهد. در این پارادایم، محدودیت‌های ساختاری تنک‌سازی (نظیر الگوهای سخت‌افزاری \lr{2:4 Structured Sparsity}) و محدودیت‌های کوانتیدگی یکنواخت (نظیر کوانتیزاسیون ۴بیتی به ازای هر کانال) به طور هم‌زمان در قالب متغیرهای کمکی فرمول‌بندی شده و تعادل بهینه‌ای میان حفظ دقت، کاهش فضای دیسک و شتاب‌بخشی استنتاج روی شتاب‌دهنده‌های مدرن تنسور کُر \lr{(Tensor Cores)} به دست می‌آید.
